\documentclass[11pt,a4paper]{article}
\usepackage[T1]{fontenc}
\usepackage{amssymb}
\usepackage{isabelle,isabellesym}
\usepackage{pdfsetup}
\hypersetup{hidelinks}
\urlstyle{rm}
\isabellestyle{it}
\newcommand{\theoremname}[1]{{\urlstyle{tt}\nolinkurl{#1}}}
\begin{document}
\pagenumbering{roman}
\title{Set-Coded Turing Machines in Isabelle/ZF}
\author{Tang Ziyi}
\maketitle

\begin{abstract}
This library defines deterministic Turing machines as finite lists of
set-coded instructions in Isabelle/ZF. It provides execution semantics,
finite-run witnesses, and an observational equivalence for finite tape
representations that differ only in stored blank cells. An explicit
first-order set-theoretic formula expresses halting. Its satisfaction in
every transitive ZF set model is equivalent to external halting of the
given machine on the given input. Finite-object names close the formula.
An injective formula code and primitive-recursive constructors certify
the numerical codes of the resulting self-input halting sentences.
Worked examples connect program execution, tape observations, and
internal satisfaction.
\end{abstract}

\tableofcontents
\clearpage
\pagenumbering{arabic}

\section{Machine interface}
A machine is a finite list of instructions over the symbols 0 and 1.
Each instruction contains an action and a next state. The actions write
0, write 1, move left, move right, or leave the tape unchanged. State 1
is initial and state 0 is final. The two entries for a positive state $q$
occur at indices $2(q-1)$ and $2(q-1)+1$. Missing entries enter the final
state without changing the tape. The final state is absorbing.
This design follows the Isabelle/HOL AFP development
\emph{Universal Turing Machine}~\cite{afp-utm}; the objects and proofs here
are developed in Isabelle/ZF.

A configuration contains a state and two finite symbol lists. The left
list starts immediately left of the head, while the right list starts
at the scanned cell. Unrepresented cells are blank. The functions
\theoremname{tm_step} and \theoremname{steps} specify one step and a
finite iteration. The predicate \theoremname{halts_on} includes the
well-formedness of the machine and input and the existence of a finite
iteration reaching state 0. Theorem \theoremname{halts_on_iff_finite_run}
gives the equivalent finite-sequence witness form.

\theoremname{Turing_Examples} defines a two-instruction machine that
scans a unary input and writes one at the first blank. The theorem
\theoremname{scan_append_result} specifies its final configuration after
$n+1$ steps. The head remains on the newly written cell. The example
proves a loop invariant, derives halting through \theoremname{halts_onI},
and uses \theoremname{steps_config_eq} to transfer the result to an input
with an extra stored blank. It provides a client of the execution and
tape interfaces without importing model theory or numerical coding.

The relations \theoremname{half_tape_eq}, \theoremname{tape_eq}, and
\theoremname{config_eq} compare observations at fixed positions relative
to the head. They ignore trailing stored blanks but do not identify
arbitrary shifts of the head. Raw list equality and observational
equality have separate uses. In particular, literal storage length is
not an output specification invariant under tape equivalence.

\section{Halting formula interface}
For a machine $P$, finite input $x$, and transitive ZF set model $M$,
the explicit formula satisfies
\[
  (M,[P,x])\models\mathit{halt\_fm}
  \quad\Longleftrightarrow\quad
  \mathrm{halts\_on}(P,x).
\]
The satisfaction theorem is
\theoremname{transitive_model_halt_fm_iff}.
The closed sentence $\mathit{halt\_sentence}(P,x)$ names both finite
parameters. The theorem \theoremname{sats_halt_sentence_iff} establishes
its corresponding satisfaction equation. Both results require
$P\in\mathit{tm\_machine}$ and $x\in\mathit{list}(\mathit{symbol})$.
They do not assert that the formula recognises arbitrary sets as valid
machines or inputs.

The proofs hold in the weaker locale \texttt{M\_ZF1\_trans}, which
assumes transitivity, Zermelo set theory, and four specified replacement
instances. The qualified satisfaction theorems can be used directly in
extensions of that locale. The unqualified theorems are the full ZF
corollaries. No existence of a transitive model is asserted.
\theoremname{Turing_Halting_Examples} derives internal truth of the
example's halting sentence using this interface and its operational
proof, without unfolding the internal formula.

Closure and satisfaction use the AFP model infrastructure
\cite{afp-models}. The \texttt{Independence\_CH} entry supplies the
locale and full-ZF interpretation in
\theoremname{ZF_Trans_Interpretations}. The packaged session therefore
retains this entry as a dependency~\cite{afp-ch}.
No CH independence theorem or choice assumption is used by the halting
interface. The underlying absoluteness method is classical
\cite[Chapter~12]{jech}; this library supplies the explicit formula and
its verified connection to the machine semantics.

\section{Effective codes}
\theoremname{Turing_Coding} assigns natural-number codes to machines and
configurations. Decoding an encoded machine recovers it. Decoding an
arbitrary natural number normalises invalid instruction fields, so the
decoder is total but not injective on all natural numbers.
\theoremname{Formula_Coding} supplies an injective code for the
set-theoretic formulas of Isabelle/ZF~\cite{paulson-zf}.

The numerical function \theoremname{pr_halt_sentence_code} satisfies
\[
  f(e)=\ulcorner\mathit{halt\_sentence}
    (\mathit{decode\_machine}(e),\mathit{numeral\_input}(e))\urcorner,
\]
where $f(e)$ denotes its application to the singleton list $[e]$.
Its membership in \texttt{prim\_rec} and its application equation are
proved in \theoremname{Turing_Halting_Codes}. This certificate uses the
bundled Isabelle/ZF primitive-recursive function class~\cite{szasz}.
It is a certificate for self-input sentence construction. A general
compiler implementing primitive-recursive functions in this machine
model is outside the library's interface.

\input{session}

\clearpage
\begin{thebibliography}{9}

\bibitem{afp-ch}
Emmanuel Gunther, Miguel Pagano, Pedro S\'anchez Terraf, and Mat\'ias Steinberg.
\emph{The Independence of the Continuum Hypothesis in Isabelle/ZF}.
Archive of Formal Proofs, 2022.
\url{https://isa-afp.org/entries/Independence_CH.html}.

\bibitem{afp-utm}
Jian Xu, Xingyuan Zhang, Christian Urban, Sebastiaan J. C. Joosten, and
Franz Regensburger.
\emph{Universal Turing Machine}.
Archive of Formal Proofs, 2019, revised 2022.
\url{https://isa-afp.org/entries/Universal_Turing_Machine.html}.

\bibitem{afp-models}
Emmanuel Gunther, Miguel Pagano, Pedro S\'anchez Terraf, and Mat\'ias Steinberg.
\emph{Transitive Models of Fragments of ZFC}.
Archive of Formal Proofs, 2022.
\url{https://isa-afp.org/entries/Transitive_Models.html}.

\bibitem{jech}
Thomas Jech.
\emph{Set Theory: The Third Millennium Edition, Revised and Expanded}.
Springer, 2003.
\url{https://doi.org/10.1007/3-540-44761-X}.

\bibitem{paulson-zf}
Lawrence C. Paulson.
\emph{The Relative Consistency of the Axiom of Choice, Mechanized Using
Isabelle/ZF}.
LMS Journal of Computation and Mathematics, 6, 198--248, 2003.

\bibitem{szasz}
Nora Szasz.
\emph{A Machine Checked Proof that Ackermann's Function is not Primitive
Recursive}.
In G\'erard Huet and Gordon Plotkin, editors, \emph{Logical Environments},
317--338. Cambridge University Press, 1993.

\end{thebibliography}

\end{document}
